% Lecture 02: If Planck Prosecuted Cagliostro for Gacha Fraud
% Timeline Prequel - Lecture 00
\documentclass[12pt]{article}
\usepackage{amsmath,amssymb,physics,bm}
\usepackage{graphicx}
\usepackage[dvipsnames]{xcolor}
\usepackage{tcolorbox}
\usepackage[colorlinks=true,
            linkcolor=MidnightBlue,
            urlcolor=RoyalBlue,
            citecolor=OliveGreen]{hyperref}

\title{Lecture 03: \\ Wigner's Friend and the Jusenkyo Trial}
\author{Written by Junhu Park}
\date{Lecture 03 \\ Ranma 1/2 Learns Quantum Mechanics}

\begin{document}

\maketitle

\begin{tcolorbox}[title=Note to the Reader]
References are intentionally omitted from this lecture for readability. For historical physicists, quantum foundations concepts, and fictional characters appearing in this story, readers are encouraged to consult Wikipedia or standard textbooks.
\end{tcolorbox}

\section*{Cast of Characters}

\begin{itemize}
    \item Ryoga Hibiki --- Plaintiff, victim of Jusenkyo's cursed spring
    \item Jusenkyo Tourism Development Corporation Manager --- Defendant
    \item Ranma Saotome --- Witness, martial artist and friend
    \item Akane Tendo --- Witness, Ranma's fiancée and observer
    \item Genma Saotome (Panda) --- Ryoga's companion, transformed guardian
    \item Professor Wigner --- Quantum theorist and expert witness
    \item Judge --- Presiding over the trial
    \item Plaintiff Counsel --- Representing Ryoga
    \item Defense Counsel --- Representing Jusenkyo Tourism Development Corporation
    \item Court Clerk --- Administrative officer of the court
\end{itemize}

\section*{Act 1. Court in Session}

\begin{tcolorbox}[title=Judge:]
The court is now in session. Ryoga Hibiki, you are called to testify. Please state your case.
\end{tcolorbox}

\begin{tcolorbox}[title=Ryoga:]
Your Honor, I fell into the Spring of Drowned Piglet. It is cursed, and since that day, I have been transformed.
\end{tcolorbox}

\begin{tcolorbox}[title={Ryoga (continued):}]
I became a pig. Not figuratively — literally a piglet named P-chan!
\end{tcolorbox}

\begin{tcolorbox}[title={Court Clerk (scribbling):}]
Plaintiff. Became Pig.
\end{tcolorbox}

\begin{tcolorbox}[title={Manager (objecting):}]
Objection, Your Honor! The plaintiff’s claim that he “became a pig” is unproven. We have no evidence of this so-called “pighood.” Mere allegation is insufficient.
\end{tcolorbox}

\begin{tcolorbox}[title=Ryoga:]
But Your Honor, isn’t my own testimony evidence? I am the one who experienced it firsthand!
\end{tcolorbox}

\begin{tcolorbox}[title=Manager:]
Self-interested testimony is not reliable evidence. The plaintiff stands to gain from this claim. We ask the court to dismiss such biased statements.
\end{tcolorbox}

\begin{tcolorbox}[title=Ryoga:]
Your Honor, how can you dismiss my own words about my condition? Am I to be silenced by my adversary’s skepticism?
\end{tcolorbox}

\begin{tcolorbox}[title=Manager:]
Skepticism is the foundation of justice. Testimony must be corroborated by independent evidence, not just self-assertion.
\end{tcolorbox}

\begin{tcolorbox}[title={Ryoga (exasperated):}]
Is this court insane? I stand before you, confessing a transformation, and you treat me as if I am fabricating fiction!
\end{tcolorbox}

\begin{tcolorbox}[title={Judge (firmly):}]
Order! The court will consider all testimony but requires evidence beyond self-assertion. Proceed to present such evidence.
\end{tcolorbox}

\section*{Act 2. Experimental Evidence}

\begin{tcolorbox}[title=Judge:]
Before we proceed, why is cold water necessary to demonstrate this transformation?
\end{tcolorbox}

\begin{tcolorbox}[title=Plaintiff Counsel:]
Your Honor, the curse activates only upon contact with cold water. It is the trigger for the transformation.
\end{tcolorbox}

\begin{tcolorbox}[title=Judge:]
Very well. Bailiff, please bring a tub of cold water.
\end{tcolorbox}

\begin{tcolorbox}[title={Bailiff (entering with tub):}]
Here is the cold water, Your Honor.
\end{tcolorbox}

\begin{tcolorbox}[nobeforeafter,sharp corners=south, colframe=black!50!white, colback=white, fontupper=\itshape, center title, title=Stage Direction:]
Ryoga steps carefully into the tub. A splash echoes through the courtroom. Silence falls. Slowly, Ryoga’s form begins to change — ears elongate, snout emerges, limbs shorten. The transformation completes: P-chan, the piglet, stands where Ryoga was moments ago.
\end{tcolorbox}

\begin{tcolorbox}[title={Judge (leaning forward):}]
Is this... the plaintiff?
\end{tcolorbox}

\begin{tcolorbox}[title=P-chan:]
Oink! I mean... Your Honor, I am still Ryoga inside. This form is not my true self.
\end{tcolorbox}

\begin{tcolorbox}[title={Judge (confused):}]
Could you clarify? Are you asserting that despite appearances, you remain the same individual?
\end{tcolorbox}

\begin{tcolorbox}[title=P-chan:]
Yes, exactly! The curse changes my body, but my mind, memories, and identity persist.
\end{tcolorbox}

\begin{tcolorbox}[title=Judge:]
Is there any way to verify this claim objectively?
\end{tcolorbox}

\begin{tcolorbox}[title=Court Clerk:]
Your Honor, there is no certified pig interpreter in this court. We cannot guarantee the accuracy of this testimony.
\end{tcolorbox}

\begin{tcolorbox}[title={Manager (smirking):}]
Indeed, Your Honor. This “testimony” is suspect, given the language barrier and the obvious self-interest.
\end{tcolorbox}

\begin{tcolorbox}[title={Judge (deadpan):}]
Noted. The court will take this under advisement.
\end{tcolorbox}

\section*{Act 3. The Mysterious Panda}

\begin{tcolorbox}[nobeforeafter,sharp corners=south, colframe=black!50!white, colback=white, fontupper=\itshape, center title, title=Stage Direction:]
Suddenly, the courtroom door bursts open with a dramatic creak. Genma Saotome, in his panda form, ambles in, holding two large signboards.
\end{tcolorbox}

\begin{tcolorbox}[title={Judge (startled):}]
What is the meaning of this? Why is a panda present in my courtroom?
\end{tcolorbox}

\begin{tcolorbox}[title={Genma (Panda) --- First Sign:}]
\textit{"The pig is asking for hot water."}
\end{tcolorbox}

\begin{tcolorbox}[title={Judge (raising an eyebrow):}]
A pig asking for hot water? That is... unusual.
\end{tcolorbox}

\begin{tcolorbox}[title={Genma (Panda) --- Second Sign:}]
\textit{"That is not the important issue."}
\end{tcolorbox}

\begin{tcolorbox}[title={Judge (considering):}]
Interesting. The panda suggests the focus should not be on the pig’s request, but on a deeper matter.
\end{tcolorbox}

\begin{tcolorbox}[title=Plaintiff Counsel:]
Your Honor, may we provide hot water to test the transformation?
\end{tcolorbox}

\begin{tcolorbox}[title=Judge:]
Proceed.
\end{tcolorbox}

\begin{tcolorbox}[nobeforeafter,sharp corners=south, colframe=black!50!white, colback=white, fontupper=\itshape, center title, title=Stage Direction:]
Hot water is brought forth. P-chan steps into the tub. A splash resonates. Slowly, the piglet’s form shifts, limbs elongate, features sharpen, and Ryoga reappears, human once more.
\end{tcolorbox}

\begin{tcolorbox}[title={Ryoga (standing, triumphant):}]
I have returned to my true form! This proves the curse is reversible, and I remain the same individual.
\end{tcolorbox}

\begin{tcolorbox}[title=Judge:]
This demonstration is compelling. The transformation is undeniable. The question remains: does this prove identity continuity?
\end{tcolorbox}

\section*{Act 4. The Identity Problem}

\begin{tcolorbox}[title=Defense Counsel:]
Your Honor, the existence of a pig before hot water, and a human after, does not prove they are the same entity. The court must consider the philosophical problem of identity continuity.
\end{tcolorbox}

\begin{tcolorbox}[title={Defense Counsel (continued):}]
Without admissible evidence linking the internal states of the pig and the human, the claim of identity is unsubstantiated.
\end{tcolorbox}

\begin{tcolorbox}[title=Plaintiff Counsel:]
But the memories and personality persist! Ryoga remembers everything as P-chan and vice versa. That continuity is the essence of identity.
\end{tcolorbox}

\begin{tcolorbox}[title=Defense Counsel:]
Memory alone is insufficient. We require objective, admissible proof. Philosophers debate whether identity is physical, psychological, or legal. The court must be cautious.
\end{tcolorbox}

\begin{tcolorbox}[title=Ryoga (frustrated):]
Is it not obvious? I am the same person, regardless of form!
\end{tcolorbox}

\begin{tcolorbox}[title=Defense Counsel:]
Obviousness is not evidence. The court demands admissible proof, not mere assertions.
\end{tcolorbox}

\section*{Act 5. Professor Wigner Appears}

\begin{tcolorbox}[title=Professor Wigner:]
Your Honor, allow me to elucidate the quantum analogy relevant to this case.

Consider Ryoga’s state as a quantum system undergoing transformations:

\[
|\text{Ryoga}\rangle
\rightarrow
|P\text{-chan}\rangle
\rightarrow
|\text{Ryoga}\rangle
\]

More generally, the system can be in a superposition:

\[
|\Psi\rangle
=
\alpha|\text{Ryoga}\rangle
+
\beta|P\text{-chan}\rangle
\]

where \(\alpha\) and \(\beta\) are complex amplitudes.

From the court’s external perspective, the combined state including observation is:

\[
|\Psi_{\text{court}}\rangle
=
\alpha|\text{Pig Observed}\rangle
+
\beta|\text{Human Observed}\rangle
\]

This illustrates observer-dependent facts: what is true depends on the observer’s knowledge.

Genma Panda acts as Wigner’s Friend — an internal observer who retains continuous information about Ryoga’s state, while the court observes only the collapsed outcome.

Measurement records, epistemic access, and observer perspectives are crucial. The court, as an external observer, lacks direct access to Ryoga’s internal state, represented by Panda’s knowledge.
\end{tcolorbox}

\begin{tcolorbox}[title=Judge:]
So, Panda is the keeper of internal information, analogous to Wigner’s Friend observing the system inside the lab?
\end{tcolorbox}

\begin{tcolorbox}[title=Professor Wigner:]
Precisely, Your Honor.
\end{tcolorbox}

\section*{Act 6. Ranma Testifies}

\begin{tcolorbox}[title=Ranma:]
I have known Ryoga for years. When he becomes P-chan, I see the same personality, memories, and fighting spirit. To me, P-chan is clearly Ryoga.
\end{tcolorbox}

\begin{tcolorbox}[title={Ryoga (beaming):}]
Finally! Someone who understands!
\end{tcolorbox}

\begin{tcolorbox}[title=Judge:]
Ranma’s testimony is noted. The court is inclined to accept this identification.
\end{tcolorbox}

\begin{tcolorbox}[title={Defense Counsel (whispering):}]
This could be problematic.
\end{tcolorbox}

\section*{Act 7. The Final Witness}

\begin{tcolorbox}[title=Judge:]
Next, we call Akane Tendo to the stand.
\end{tcolorbox}

\begin{tcolorbox}[title={Ryoga (nervous):}]
Akane? Oh no...
\end{tcolorbox}

\begin{tcolorbox}[title=Akane:]
I see P-chan as just a cute pig. There is no human identity remaining in that form.
\end{tcolorbox}

\begin{tcolorbox}[title={Ranma (alarmed):}]
Wait... if Akane says that, then... the social implications...
\end{tcolorbox}

\begin{tcolorbox}[title={Ranma (to Judge):}]
Your Honor, I must formally retract all my previous testimony.
\end{tcolorbox}

\begin{tcolorbox}[title={Ryoga (angrily):}]
Traitor! How can you deny me now?
\end{tcolorbox}

\begin{tcolorbox}[title={Ranma (frustrated):}]
Because if I confirm P-chan is Ryoga, it would ruin him socially — especially with Akane involved. It’s a disaster waiting to happen.
\end{tcolorbox}

\begin{tcolorbox}[title=Ryoga:]
I don’t get it. What disaster?
\end{tcolorbox}

\begin{tcolorbox}[title=Ranma:]
It’s complicated, but trust me — silence is the only option.
\end{tcolorbox}

\section*{Act 8. Social Collapse Avoidance}

\begin{tcolorbox}[title=Akane:]
P-chan is a cute pig. That’s all. Ryoga is Ryoga. They are not the same.
\end{tcolorbox}

\begin{tcolorbox}[title={Akane (continued):}]
I see Ryoga as a human friend, and P-chan as a separate, adorable animal.
\end{tcolorbox}

\begin{tcolorbox}[title={Judge (looking around):}]
Why does the courtroom seem so relieved now?
\end{tcolorbox}

\begin{tcolorbox}[title=Professor Wigner:]
This is what I call “Social Collapse Avoidance Success.”

If Ranma publicly confirmed

\[
\text{Observation}
:
P\text{-chan}=\text{Ryoga}
\]

it would cause a

\[
\Rightarrow
\text{Catastrophic Social Wavefunction Collapse}
\]

destroying Ryoga’s personal life and relationships.

By denying the identity, Akane preserves social harmony and prevents this collapse.
\end{tcolorbox}

\begin{tcolorbox}[title=Ranma:]
In quantum terms, the measurement outcome is observer-dependent, and social consequences are the “wavefunction collapse” of reputation.
\end{tcolorbox}

\section*{Act 9. Final Verdict}

\begin{tcolorbox}[title=Judge:]
Given the inability to legally establish identity between Ryoga and P-chan, the court dismisses the claim.
\end{tcolorbox}

\begin{tcolorbox}[title=Ryoga:]
I accept defeat, though it pains me. At least Akane sees P-chan as just a pig.
\end{tcolorbox}

\begin{tcolorbox}[title={Judge (confused):}]
But the damages were enormous. Why dismiss?
\end{tcolorbox}

\begin{tcolorbox}[title={Akane (calling out affectionately):}]
P-chan!
\end{tcolorbox}

\begin{tcolorbox}[title={Ryoga (smiling faintly):}]
That’s a small mercy.
\end{tcolorbox}

\begin{tcolorbox}[title={Judge (finally understanding):}]
I see. Sometimes, the law must bow to social reality.
\end{tcolorbox}

\section*{Epilogue: Wigner's Friend}

\begin{tcolorbox}[title=Professor Wigner:]
This trial serves as a metaphor for the Wigner’s Friend thought experiment, highlighting the subtleties of observer-relative reality.

Key points:

\begin{itemize}
    \item Ryoga knows the truth of his own identity — the internal observer.
    \item The court lacks access to this internal state — external observer.
    \item Genma Panda holds the internal state information, acting as Wigner’s Friend.
    \item Akane is the dangerous observer whose measurement collapses the identity superposition socially.
    \item Ranma updates his testimony understanding the social consequences of observation.
\end{itemize}

Mathematically, the information held by Panda exceeds that of the court:

\[
I_{\text{Panda}} > I_{\text{Court}}
\]

Observers Akane and Ranma have different observables:

\[
\mathcal O_{\text{Akane}}
\neq
\mathcal O_{\text{Ranma}}
\]

The knowledge gap between observers is nonzero:

\[
\Delta K_{\text{observer}}
\neq 0
\]

This reflects the epistemic limitations and the role of measurement in quantum mechanics and social systems.

In law and physics alike, identity and reality depend on the observer’s perspective and information.

\[
\boxed{\text{Optimal Defeat}}
\]

\noindent
Sometimes, the best outcome is to accept uncertainty and preserve harmony.
\end{tcolorbox}

\begin{tcolorbox}[title={Genma (Panda):}]
So why was I here?
\end{tcolorbox}

\begin{tcolorbox}[title=Ranma:]
Good question.
\end{tcolorbox}

\begin{tcolorbox}[title=Ryoga:]
No idea.
\end{tcolorbox}

\begin{tcolorbox}[title={Akane (smiling):}]
P-chan is adorable!
\end{tcolorbox}


\section*{Further Reading}

For quantum mechanics, quantum information, and quantum foundations, consult standard textbooks.

For Eugene Wigner, Wigner's Friend, Ranma 1/2, Ryoga Hibiki, Akane Tendo, Genma Saotome, and related cultural references, Wikipedia provides sufficient background for the intended audience of this lecture series.

\end{document}
